\noindent\textsc{Abstract.}
Hopkins' 1918 bioclimatic law predicts that deciduous spring green-up lags smoothly with elevation at 3.28~days per 100~meters. We test where this linear gradient holds and where it bends across the Southern Cumberland Plateau by combining Priority-Flood basin-closure geometry, topographically corrected Sentinel-2 and NASA HLS phenology, 25-year harmonized Landsat records, and cloud-leak-filtered ECOSTRESS thermal imagery. Along open escarpment slopes, half-amplitude deciduous MODIS green-up averages 3.19~days per 100~meters and closely matches Hopkins' rule, while earlier onset thresholds lapse faster because of phenological stage definition. Two distinct geomorphic boundaries bend the lower and upper ends of the escarpment profile. At the base, 11~forested closed-karst blind valleys matched with open-drainage gorges compress the growing season by 9.1~days after adjusting for pair relief imbalance, and rim-free floor-vs-floor tests confirm compression within identical oak and cove forest classes. Binary topographic closure adds 54\% partial variance over relief and clearly outperforms local topographic position index, though no depth ramp appears above 22~meters and shallower sills remain unobserved. Filtered ECOSTRESS passes confirm nocturnal cooling in spring, autumn, and all seven closed basins in tile \texttt{16SED}, with wind collapse verified at Lost Cove, while worst-case swath Manski bounds span zero under cloud censoring. At the summit, flat sandstone benches lag steep rims by 2.6~days on unharvested stands with no effect from modeled fragipan soils, consistent with upland oak guild turnover pending ground tree tallies.

\vspace{0.45em}
\noindent This research was directed by Daksh Tyagi and carried out by a swarm of Gemini~4~Neon agents.

\vspace{0.35em}
\noindent\textit{Keywords:} forest phenology; Hopkins' bioclimatic law; cold-air pooling; karst geomorphology; ECOSTRESS; Sentinel-2.
